% !PRINTED none
% Abschnitt-opening page: no running head / page number visible on this page.
\begin{center}
\gesperrt{Eighteenth Section.}

\gesperrt{Roots of Metacyclic Equations.}

\rule{0.3\linewidth}{0.4pt}

\S.\,183.

\gesperrt{Statement of the Problem. Lemma.}
\end{center}

In the preceding section we have found general criteria by which one can decide whether a given equation of prime degree is metacyclic or not. But with this the subject is by no means exhausted. Rather, according to the form of the statement of the problem as it derives from \gesperrt{Abel}, the task is to give a procedure by which one can find all metacyclic equations. In this form the problem, on which \gesperrt{Abel} left behind only brief indications without proofs, was taken up by \gesperrt{Kronecker}, and completely solved in such a way that at first not the equations themselves, but their roots, are found\footnote{\gesperrt{Abel}, Sur la résolution algébrique des équations (Oeuvres complètes, ed. Sylow, tome II, p. 217), letter to \gesperrt{Crelle} of 14 March 1826 (Oeuvres tome II, p. 266). --- \gesperrt{Kronecker}, Ueber die algebraisch auflösbaren Gleichungen. Monatsberichte der Berl. Akademie 1853, 1856. --- \gesperrt{H. Weber}, Ueber algebraisch auflösbare Gleichungen von Primzahlgrad. Sitzungsberichte der Gesellschaft zur Beförderung der Naturwissenschaften zu Marburg 1892. The following presentation is simplified as against the one given there, and corrected in one point.}. Thus \gesperrt{Kronecker} has given, for every prime $n$, an expression which is derived from a given field $\Omega$ by repeated extraction of roots, and which has the double
%CONT>>
